\documentclass[11pt]{article}
\usepackage[margin=1in]{geometry}
\usepackage{amsmath,amssymb,amsthm,mathtools}
\usepackage{enumitem}
\usepackage{xcolor}
\usepackage{booktabs}
\usepackage[colorlinks=true,linkcolor=blue!50!black,citecolor=blue!50!black,urlcolor=blue!50!black]{hyperref}

\newtheorem{theorem}{Theorem}
\newtheorem{lemma}[theorem]{Lemma}
\newtheorem{proposition}[theorem]{Proposition}
\newtheorem{corollary}[theorem]{Corollary}
\theoremstyle{definition}
\newtheorem{assumption}{Assumption}
\newtheorem{definition}[theorem]{Definition}
\newtheorem{remark}[theorem]{Remark}

\newcommand{\E}{\mathbb{E}}
\newcommand{\Prob}{\mathbb{P}}
\newcommand{\R}{\mathbb{R}}
\newcommand{\Nb}{N_{\mathrm b}}
\newcommand{\Xmax}{X_{\max}}
\newcommand{\eps}{\varepsilon}
\newcommand{\cE}{\mathcal{E}}
\newcommand{\cG}{\mathcal{G}}
\newcommand{\ind}{\mathbf 1}
\newcommand{\UCB}{\mathrm{UCB}}
\DeclareMathOperator*{\argmax}{arg\,max}

\title{Dec-ZoomSIB-Relay without a server: proofs\\[2pt]
\large Lower bound, regret, communication, robustness and agreement}
\author{CS6007 project notes}
\date{October 2026}

\begin{document}
\maketitle

\begin{abstract}
We analyse the serverless version of Fed-ZoomSIB, where $N$ agents communicate only along the edges of a graph
$\cG$. Every message is flooded, so it reaches every agent within $\Lambda$ rounds; on a static connected graph
$\Lambda$ is the diameter $D$. We prove five results. (1) No serverless algorithm beats $\Omega((NT)^{2/3})$
(Theorem~\ref{thm:lower}). (2) Dec-ZoomSIB-Relay has regret equal to the server bound plus an additive
$O(L_f\,N\Lambda\,\Nb)$, that is, $\tilde O\big((d NT)^{2/3}+N\Lambda(NT)^{1/3}\big)$ (Theorem~\ref{thm:regret}).
(3) It sends $O\big(|E|\,N\,\Nb^2\log_{1+\gamma}(NT)\big)$ scalars, at most logarithmic in $T$ (Theorem~\ref{thm:comm}).
(4) Time-varying, unreliable and directed graphs only change $\Lambda$ (Corollary~\ref{cor:robust}). (5) If agents
freeze different directions, the regret grows by at most $4L_{f'}\Xmax\,\eta\, NT$, where $\eta$ is the worst
direction error, and the cost is zero when Phase~1 pools exactly (Theorem~\ref{thm:agree}). The analysis uses
\emph{synchronised epochs}: all agents act on the same delayed pooled table plus their own recent pulls. This
keeps the confidence argument of the server case unchanged and removes the $\sqrt N$ loss of the earlier sketch.
\end{abstract}

\tableofcontents

%======================================================================
\section{Setting, notation and imported results}\label{sec:setting}
%======================================================================

\paragraph{Bandit model.} There are $N$ agents and $T$ rounds. At round $t$ agent $i$ observes $K$ contexts
$X_{i,t}=(x_{i,t,1},\dots,x_{i,t,K})$, drawn i.i.d.\ from $\mathcal D$ on $\R^d$, independently across agents
and rounds. It pulls an arm $a_{i,t}$ and observes
\[
y_{i,t}=f(\langle x_{i,t,a_{i,t}},\theta^*\rangle)+\eta_{i,t}.
\]
Both $\theta^*$ (with $\|\theta^*\|_1=1$) and $f$ are unknown and shared by all agents. Write $z=\langle x,\theta^*\rangle$
for the \emph{true index} of an arm $x$. The network regret is
\[
R_T(N)=\sum_{i=1}^N\sum_{t=1}^T r_{i,t},\qquad
r_{i,t}=\max_{a\le K} f(\langle x_{i,t,a},\theta^*\rangle)-f(\langle x_{i,t,a_{i,t}},\theta^*\rangle).
\]

\paragraph{Order of pulls.} Within a round, all $N$ agents act in parallel. For the analysis we order all pulls
\emph{lexicographically} by (round, agent): $(t,i)\prec(t',i')$ iff $t<t'$, or $t=t'$ and $i<i'$. Let
$\mathcal F_{t,i}$ be the $\sigma$-field generated by all contexts, actions and rewards strictly before $(t,i)$ in
this order, together with $X_{i,t}$ and the algorithm's internal randomness.

\begin{assumption}[DBG Assumptions 1--3]\label{ass}
\begin{enumerate}[nosep,label=(\roman*)]
\item $\eta_{i,t}$ is independent of the contexts and conditionally $\sigma$-sub-Gaussian:
$\E[e^{\lambda\eta_{i,t}}\mid\mathcal F_{t,i},a_{i,t}]\le e^{\lambda^2\sigma^2/2}$ for all $\lambda\in\R$.
\item $\E[S_j(X)^2]\le M$ for the score $S=-\nabla\log p$.
\item Each coordinate of $X$ is $\kappa$-sub-Gaussian. On $|z|\le\Xmax$ we have $|f|\le L_f$ and $|f'|\le L_{f'}$, where
$\Xmax=\kappa\sqrt{2\log(2dNTK/\delta)}$. With probability $\ge1-\delta$, every one of the $NTK$ contexts satisfies
$\|x\|_\infty\le\Xmax$. Call this event $E_x$.
\item $\mu^*=\E[f'(\langle X,\theta^*\rangle)]>0$.
\end{enumerate}
\end{assumption}

\paragraph{Communication model.} The agents are the vertices of a graph $\cG=(V,E)$ with $|V|=N$. Each round has
two steps: a \emph{decision step}, in which every agent pulls an arm, and then a \emph{communication step}, in which
every agent may send messages to its neighbours. A message sent in the communication step of round $s$ is
received before the decision step of round $s+1$. \emph{Flooding} means that an agent forwards every message the
first time it receives it, to all its neighbours, and never again.

\begin{assumption}[Delivery time]\label{ass:delivery}
There is a known integer $\Lambda\ge1$ such that a message that is flooded from any agent in the communication step
of round $s$ has reached every agent by the end of the communication step of round $s+\Lambda-1$. Each agent knows
$N$, $\Lambda$ and the round number.
\end{assumption}

\begin{lemma}[Flooding on a static graph]\label{lem:flood}
If $\cG$ is static, connected and has diameter $D$, Assumption~\ref{ass:delivery} holds with $\Lambda=D$.
\end{lemma}
\begin{proof}
Let agent $a$ flood a message in step $s$. We show by induction on $h\ge1$ that every agent at distance $h$ from
$a$ holds the message after step $s+h-1$. For $h=1$: neighbours of $a$ receive it in step $s$. For the inductive step,
let $v$ be at distance $h+1$. It has a neighbour $u$ at distance $h$, which holds the message after step $s+h-1$. If
$u$ received it in step $s+h-1$ or earlier, it forwards it no later than step $s+h$. So $v$ holds it after step $s+h$.
Every agent is at distance at most $D$.
\end{proof}

\paragraph{Imported results.} The following results are proved elsewhere and used as black boxes.
Here DBG refers to Dey, Bhore \& Ghosh (2026).

\begin{lemma}[Stein estimator; DBG Lemma B.1 with Lemma 2.1]\label{lem:stein}
Let $\hat\theta=\frac1n\sum_{s=1}^n\varphi_\tau(y_sS(x_s))$ be the truncated Stein estimator on $n$ i.i.d.\ pairs, with
$\tau=\sqrt{3(\sigma^2+L_f^2)Mn/\log(2d/\delta)}$, and let $\hat\theta_0=\hat\theta/\|\hat\theta\|_1$. If
$C_\theta d\sqrt{\log(2d/\delta)/n}\le\mu^*/2$, then with probability $\ge1-\delta$,
\[
\|\hat\theta_0-\theta^*\|_1\le\eps_0(n):=\frac{4C_\theta d}{\mu^*}\sqrt{\frac{\log(2d/\delta)}{n}},
\qquad C_\theta=\tfrac{4\sqrt3}{3}+2\sqrt2\sqrt{M(\sigma^2+L_f^2)}.
\]
\end{lemma}

\begin{lemma}[Single-agent lower bound; DBG Theorem 5.1]\label{lem:dbglower}
There are $c>0$ and $n_0$ such that for every horizon $n\ge n_0$ and every single-agent policy $\pi$,
$\sup\E[R_n(\pi)]\ge c\,n^{2/3}$. The supremum is over instances that satisfy the regularity assumptions (with
$d=1$, Gaussian contexts, and a number of arms that may grow logarithmically with $n$).
\end{lemma}

\begin{lemma}[Martingale concentration; Azuma--Hoeffding for sub-Gaussian increments]\label{lem:azuma}
Let $(\xi_m)_{m\ge1}$ be adapted to a filtration $(\mathcal H_m)$ and satisfy
$\E[e^{\lambda\xi_m}\mid\mathcal H_{m-1}]\le e^{\lambda^2\sigma^2/2}$ for every $\lambda\in\R$. Then for every fixed $p\ge1$ and
$\delta'\in(0,1)$,
\[
\Prob\Big(\Big|\sum_{m=1}^p\xi_m\Big|>\sigma\sqrt{2p\log(2/\delta')}\Big)\le\delta'.
\]
\end{lemma}

%======================================================================
\section{Lower bound}\label{sec:lower}
%======================================================================

\begin{theorem}[No serverless algorithm beats $(NT)^{2/3}$]\label{thm:lower}
Let $\mathcal A$ be any $N$-agent algorithm, on any graph and with any communication. This includes the extreme case
where every agent sees every other agent's data after every round. Then for $NT\ge n_0$,
\[
\sup_{\text{instances}}\E\big[R_T(N;\mathcal A)\big]\ \ge\ c\,(NT)^{2/3}.
\]
\end{theorem}

\begin{proof}
The proof is a reduction: from $\mathcal A$ we build a single-agent policy $\pi$ with horizon $n=NT$ whose regret
has the same law as the network regret of $\mathcal A$. Then we apply Lemma~\ref{lem:dbglower}.

\emph{Step 1: relabelling time.} Map the single-agent round $u\in\{1,\dots,NT\}$ to the pair
\[
u\ \longleftrightarrow\ (t,i),\qquad u=(t-1)N+i,\quad t\in[T],\ i\in[N].
\]

\emph{Step 2: the policy $\pi$.} Policy $\pi$ stores the complete state of $\mathcal A$: every agent, every
message buffer, and every random seed. It draws the seeds itself. At its round $u=(t,i)$:
\begin{enumerate}[nosep,label=(\alph*)]
\item it receives the context set $X_u$ and gives it to the simulated agent $i$;
\item the simulated agent $i$ chooses $a_u$ from its state after round $t-1$, using exactly the rule of $\mathcal A$;
\item $\pi$ plays $a_u$ and stores the reward $y_u$ without giving it to the simulation;
\item if $i=N$, $\pi$ gives the $N$ stored rewards of round $t$ to the simulated agents and runs one communication
step of $\mathcal A$.
\end{enumerate}
Each action $a_u$ depends only on $X_u$, on rewards of rounds $u'\le(t-1)N<u$, and on $\pi$'s own randomness.
So $\pi$ is a valid single-agent policy.

\emph{Step 3: the laws coincide.} Let $H_t$ be the history (contexts, actions, rewards) of rounds $1,\dots,t$ of
$\mathcal A$, and $H'_t$ the history of rounds $1,\dots,tN$ of $\pi$, grouped in blocks of $N$. We show by induction
on $t$ that $H'_t\overset{d}{=}H_t$. The case $t=0$ is trivial. Assume $H'_{t-1}\overset{d}{=}H_{t-1}$.
\begin{itemize}[nosep]
\item In both systems the $N$ context sets of round $t$ are i.i.d.\ from $\mathcal D^K$ and independent of the past.
\item Given the past and the contexts, the $N$ actions are the same measurable function of these and of the seeds,
in both systems. This is because $\pi$ withholds the round-$t$ rewards.
\item Given the actions, each reward is $f(\langle x,\theta^*\rangle)+\eta$, with noise of the same conditional law in both systems.
\end{itemize}
So $H'_t\overset{d}{=}H_t$. The regret is the same function of the history in both systems, namely the sum of the
$NT$ per-pull gaps. Hence $R_{NT}(\pi)\overset{d}{=}R_T(N;\mathcal A)$ on every instance, and
\[
\sup\E R_T(N;\mathcal A)=\sup\E R_{NT}(\pi)\ \ge\ c\,(NT)^{2/3}
\]
by Lemma~\ref{lem:dbglower}.
\end{proof}

\begin{remark}
The proof uses only two facts: the agents are statistically identical (they share $\theta^*$, $f$ and $\mathcal D$),
and within a round they act in parallel. It does not depend on $\cG$. A lower bound that grows with the graph
(an additive $\Omega(\cdot\,\Lambda)$ term) would need a different argument; this is left open.
\end{remark}

%======================================================================
\section{The algorithm as analysed}\label{sec:alg}
%======================================================================

The algorithm has a grid of bins on the index line and two phases. The table below fixes the notation.

\begin{center}
\begin{tabular}{@{}lp{0.62\textwidth}@{}}
\toprule
Symbol & Meaning \\ \midrule
$T_0$ & exploration rounds per agent; Phase~2 starts at round $T_1+1$, $T_1=T_0+\Lambda$\\
$W=\Xmax$, $\Delta$ & window and bin width;\newline bins $I_j=(-W+(j-1)\Delta,\,-W+j\Delta]$, $j\le\Nb=\lceil2W/\Delta\rceil$\\
$L=2\Lambda$ & slot length; slot boundaries $\beta_m=T_1+mL$, $m\ge0$\\
$\cE$ & set of epochs (a subset of slot boundaries); $\beta_0\in\cE$\\
$O_{k,j}(s)$ & number of Phase-2 pulls by agent $k$ into bin $j$ in rounds $\le s$\\
$G_j(s)=\sum_kO_{k,j}(s)$ & network count of bin $j$ up to round $s$ (similarly the reward sum $G^S_j(s)$)\\
$e(t)$ & reference epoch at round $t$, defined in \eqref{eq:ref}\\
\bottomrule
\end{tabular}
\end{center}

\paragraph{Phase 1 (rounds $1,\dots,T_1$).} Every agent pulls uniformly at random in rounds $1,\dots,T_1$. After round
$T_0$, agent $k$ floods the record
\[
b_k=\sum_{t\le T_0}\varphi_\tau\big(y_{k,t}S(x_{k,t,a_{k,t}})\big)\in\R^d,
\]
with $\tau$ set for the pooled size $n=NT_0$. By Assumption~\ref{ass:delivery}, every agent holds all $N$ records
before round $T_1+1$. Every agent then computes the same direction
\[
\hat\theta_0=\frac{\sum_kb_k}{\|\sum_kb_k\|_1}.
\]
Pulls in rounds $T_0<t\le T_1$ are uniform and are not used.

\paragraph{Bins.} An arm $x$ belongs to bin $b(x)$, the bin with $\langle x,\hat\theta_0\rangle\in I_{b(x)}$.

\paragraph{Phase 2: synchronised epochs.} Every agent keeps a log of its own Phase-2 pulls, so it can compute
$O_{k,\cdot}(s)$ and $O^S_{k,\cdot}(s)$ for any past round $s$.
\begin{enumerate}[nosep,label=(\alph*)]
\item \emph{Trigger.} In the communication step of a slot boundary $\beta_m$ ($m\ge1$), agent $l$ raises an
\emph{alarm} if, for some bin $j$,
\begin{equation}\label{eq:trigger}
\delta_{l,j}(\beta_m):=O_{l,j}(\beta_m)-O_{l,j}(\epsilon_m)\ >\ \gamma\max\big(G_j(\epsilon'_m),1\big),
\end{equation}
where $\epsilon_m=\max\{e\in\cE:e\le\beta_{m-1}\}$ and $\epsilon'_m=\max\{e\in\cE:e\le\beta_{m-2}\}$ (with
$\beta_{-1}:=\beta_0$). In words: $l$'s own pulls into bin $j$ since the last epoch exceed $\gamma$ times the
network count it is currently acting on. The alarm (one scalar: $m$) is flooded. If at least one agent raises an alarm
at $\beta_m$, then $\beta_m\in\cE$.
\item \emph{Snapshot.} When agent $k$ first learns that $\beta_m\in\cE$, it floods its \emph{snapshot}: the bins $j$ with
$O_{k,j}(\beta_m)\neq O_{k,j}(\epsilon_m)$, each with the pair of increments, plus the epoch label.
\item \emph{Acting.} At round $t$ in slot $m+1$ (that is, $\beta_m<t\le\beta_{m+1}$), agent $i$ uses the reference epoch
\begin{equation}\label{eq:ref}
e(t)=\max\{e\in\cE:\ e\le\beta_{m-1}\},
\end{equation}
which is the same for all agents. It acts on the table
\[
n_{i,j}(t)=G_j(e(t))+\big[O_{i,j}(t-1)-O_{i,j}(e(t))\big],\qquad
S_{i,j}(t)=G^S_j(e(t))+\big[O^S_{i,j}(t-1)-O^S_{i,j}(e(t))\big].
\]
Among the bins its arms fall into, it picks $J=\argmax_j\UCB_{i,j}(t)$, where
$\UCB_{i,j}=S_{i,j}/n_{i,j}+w(n_{i,j})$ and $w(0)=\infty$, and plays any arm in bin $J$.
\end{enumerate}
The confidence width is
\[
w(n)=c_\delta\sqrt{2/n},\qquad c_\delta=\sigma\sqrt{2\log\big(4N\Nb(NT)^2/\delta\big)}.
\]

\begin{lemma}[The protocol is well defined]\label{lem:welldef}
\begin{enumerate}[nosep,label=(\roman*)]
\item In the communication step of $\beta_m$, every agent knows $\epsilon_m$, $\epsilon'_m$ and $G(\epsilon'_m)$.
\item Before round $\beta_m+1$, every agent knows $G(e)$ and $G^S(e)$ for each epoch $e\le\beta_{m-1}$.
\end{enumerate}
\end{lemma}
\begin{proof}
An alarm for $\beta_{m'}$ is flooded in step $\beta_{m'}$ and reaches everyone by step $\beta_{m'}+\Lambda-1$. Each
snapshot is flooded at the latest in step $\beta_{m'}+\Lambda$ and reaches everyone by step $\beta_{m'}+2\Lambda-1$.
Since $2\Lambda=L$, both are complete before round $\beta_{m'}+L+1=\beta_{m'+1}+1$.

For (ii), take $m'\le m-1$: everything about epoch $\beta_{m'}$ is known before round $\beta_{m'+1}+1\le\beta_m+1$.

For (i), whether $\beta_{m'}\in\cE$ is known by step $\beta_{m'}+\Lambda-1<\beta_m$ for all $m'\le m-1$, so $\epsilon_m$
and $\epsilon'_m$ are known. $G(\epsilon'_m)$ is known before round $\beta_{m-1}+1$, by (ii) applied to $m-1$.
Finally, every agent can sum snapshots to get $G(e)$, because $G_j(e)-G_j(e_{\rm prev})$ is the sum over $k$ of the increments in the snapshots for $e$.
\end{proof}

\begin{remark}[Relation to the code]\label{rem:code}
The flooding code (\texttt{src/dec/phase2/flood.py}) works differently. Each agent republishes its own table when its
\emph{own} trigger fires, and each agent acts on the freshest copy of every record it holds. That uses more information,
but it changes the confidence argument. In the code, whether a pull of agent $k$ is in agent $i$'s view depends on
when $k$ next publishes, which is decided \emph{after} the pull. Agent $i$'s view is then a union of $N$ per-agent
prefixes. Bounding each prefix separately costs a factor $\sqrt N$ in the width (this was the gap in the earlier sketch).
The epoch protocol above avoids this: every agent's view is one common prefix plus its own pulls (Lemma~\ref{lem:conf}).
E23 measures the code version. Implementing the epoch version and checking that it behaves the same is a natural next experiment.
\end{remark}

%======================================================================
\section{Regret upper bound}\label{sec:regret}
%======================================================================

\subsection{Good events and bin geometry}

Throughout this section, agent $k$ may use its own direction $\hat\theta_k$; the protocol above has
$\hat\theta_k=\hat\theta_0$ for every $k$. This generality is used again in Section~\ref{sec:agree}. Define the
\emph{projection error}
\[
\eps_{\rm p}=\Xmax\max_{k\le N}\|\hat\theta_k-\theta^*\|_1 .
\]

\begin{lemma}[Phase 1 is exactly centralised]\label{lem:exact}
$\sum_kb_k/(NT_0)$ is the truncated Stein estimator on $n=NT_0$ i.i.d.\ samples. Hence, if $NT_0$ satisfies the
sample-size condition of Lemma~\ref{lem:stein}, there is an event $E_\theta$ with $\Prob(E_\theta)\ge1-\delta$ on which
$\|\hat\theta_0-\theta^*\|_1\le\eps_0(NT_0)$.
\end{lemma}
\begin{proof}
Truncation $\varphi_\tau$ acts on each summand separately, with the same $\tau$ at every agent, so
$\sum_kb_k=\sum_{k}\sum_{t\le T_0}\varphi_\tau(y_{k,t}S(x_{k,t}))$ is the sum over the pooled sample. Each pull is
uniform over the $K$ i.i.d.\ contexts, so $(x_{k,t},y_{k,t})$ are i.i.d.\ across $(k,t)$ with the law of a
single-agent Phase-1 sample. Apply Lemma~\ref{lem:stein} with $n=NT_0$.
\end{proof}

Let $\tilde I_j=\big(I_j+[-\eps_{\rm p},\eps_{\rm p}]\big)\cap[-\Xmax,\Xmax]$ be bin $j$ widened by $\eps_{\rm p}$ on
each side, and
\[
m_j=\sup_{z\in\tilde I_j}f(z),\qquad \omega=L_{f'}(\Delta+2\eps_{\rm p}).
\]

\begin{lemma}[Bin geometry]\label{lem:geom}
On $E_x$, for every agent $k$ and every arm $x$ it sees:
\begin{enumerate}[nosep,label=(\roman*)]
\item $x$ falls in some bin, i.e.\ $|\langle x,\hat\theta_k\rangle|\le W$;
\item if $x$ is in bin $j$ under $\hat\theta_k$, then its true index $z=\langle x,\theta^*\rangle$ lies in $\tilde I_j$;
\item hence $m_j-\omega\le f(z)\le m_j$.
\end{enumerate}
\end{lemma}
\begin{proof}
(i) $|\langle x,\hat\theta_k\rangle|\le\|x\|_\infty\|\hat\theta_k\|_1=\|x\|_\infty\le\Xmax=W$.
(ii) By H\"older, $|\langle x,\hat\theta_k\rangle-z|\le\|x\|_\infty\|\hat\theta_k-\theta^*\|_1\le\eps_{\rm p}$. Also
$|z|\le\|x\|_\infty\|\theta^*\|_1\le\Xmax$.
(iii) $f(z)\le m_j$ by definition. The interval $\tilde I_j$ has length at most $\Delta+2\eps_{\rm p}$, and $f$ is
$L_{f'}$-Lipschitz on it, so $f(z)\ge m_j-\omega$.
\end{proof}

\subsection{Confidence}

For each bin $j$, list all Phase-2 pulls into bin $j$ in lexicographic order, and let $\mu_s=f(z_s)$ and $\eta_s$ be
the mean and the noise of the $s$-th such pull.

\begin{lemma}[What an agent sees]\label{lem:view}
The pulls counted in $n_{i,j}(t)$ are the union of two sets:
\begin{enumerate}[nosep,label=(\roman*)]
\item a \emph{common prefix}: all pulls into bin $j$ in rounds $\le e(t)$, i.e.\ the first $G_j(e(t))$ pulls of bin $j$'s
lexicographic list;
\item an \emph{own segment}: agent $i$'s own pulls into bin $j$ in rounds $e(t)<t'\le t-1$, i.e.\ a contiguous block of
agent $i$'s own bin-$j$ pulls.
\end{enumerate}
\end{lemma}
\begin{proof}
Immediate from the definition of $n_{i,j}(t)$ and $G_j(e)=\sum_kO_{k,j}(e)$. A pull in round $t'\le e(t)$ precedes,
in lexicographic order, every pull in a later round.
\end{proof}

\begin{lemma}[Confidence]\label{lem:conf}
There is an event $E_\eta$ with $\Prob(E_\eta)\ge1-\delta$ on which, for every agent $i$, bin $j$ and Phase-2 round $t$
with $n_{i,j}(t)\ge1$,
\[
\Big|\frac{S_{i,j}(t)}{n_{i,j}(t)}-\frac1{n_{i,j}(t)}\sum_{s\in\mathrm{seen}_{i,j}(t)}\mu_s\Big|\le w\big(n_{i,j}(t)\big),
\]
where $\mathrm{seen}_{i,j}(t)$ is the set of pulls counted in $n_{i,j}(t)$.
\end{lemma}
\begin{proof}
By Lemma~\ref{lem:view}, $S_{i,j}-\sum_{\rm seen}\mu_s=P+Q$, where $P$ is the noise sum over a prefix of length $p$ of
bin $j$'s list, and $Q$ is the noise sum over a contiguous block of length $q$ of agent $i$'s own bin-$j$ pulls, with
$p+q=n_{i,j}$.

\emph{Prefixes.} Let $\xi_m$ be the noise of the $m$-th pull into bin $j$, and $\mathcal H_m$ the information up to and
including the choice of that pull. Whether a given pull goes to bin $j$ is decided from $\mathcal F_{t,i}$, so
$(\xi_m)$ is a martingale difference sequence with $\sigma$-sub-Gaussian increments by Assumption~\ref{ass}(i). If
bin $j$ receives fewer than $m$ pulls, pad the sequence with independent $\mathcal N(0,\sigma^2)$ variables; this changes
no actual sum. Lemma~\ref{lem:azuma} applies for each fixed $p\le NT$.

\emph{Own blocks.} The same argument applies to the subsequence of agent $i$'s bin-$j$ pulls, and to the sum from its
$(a+1)$-th to its $b$-th element, for each fixed pair $0\le a<b\le T$.

\emph{Union bound.} There are at most $\Nb\cdot NT$ prefixes and $N\cdot\Nb\cdot T^2$ own blocks, in total at most
$2N\Nb(NT)^2$ events. Applying Lemma~\ref{lem:azuma} to each with $\delta'=\delta/(2N\Nb(NT)^2)$, with probability
$\ge1-\delta$ all of them satisfy
\[
|P|\le c_\delta\sqrt p,\qquad |Q|\le c_\delta\sqrt q .
\]
Since $\sqrt p+\sqrt q\le\sqrt{2(p+q)}$, we get $|P+Q|\le c_\delta\sqrt{2n_{i,j}}$. Dividing by $n_{i,j}$ gives the claim.
\end{proof}

\subsection{Regret of one pull}

Let $E=E_x\cap E_\theta\cap E_\eta$. Then $\Prob(E^c)\le3\delta$.

\begin{lemma}[Per-pull regret]\label{lem:perpull}
On $E_x\cap E_\eta$, if agent $i$ pulls from bin $J$ at Phase-2 round $t$, then
\[
r_{i,t}\le\min\big\{2L_f,\ 2w(n_{i,J}(t))+2\omega\big\}.
\]
This holds for any directions $\hat\theta_k$ (all equal or not) and for any set of available bins.
\end{lemma}
\begin{proof}
$r_{i,t}\le2L_f$ because $|f|\le L_f$ on $E_x$. For the second bound, assume $n_{i,J}\ge1$ (otherwise $w=\infty$).
Let $x^*$ be the best arm, in bin $j^*$ under $\hat\theta_i$. By Lemma~\ref{lem:geom}, $f(\langle x^*,\theta^*\rangle)\le m_{j^*}$.

\emph{The best bin looks good.} If $n_{i,j^*}=0$, then $\UCB_{j^*}=\infty$. Since $J$ maximises the UCB,
$\UCB_J=\infty$, contradicting $n_{i,J}\ge1$. So $n_{i,j^*}\ge1$. Every seen pull of bin $j^*$ has mean
$\ge m_{j^*}-\omega$ (Lemma~\ref{lem:geom}), so by Lemma~\ref{lem:conf},
\[
\UCB_{j^*}\ \ge\ (m_{j^*}-\omega)-w(n_{i,j^*})+w(n_{i,j^*})\ =\ m_{j^*}-\omega.
\]

\emph{The chosen bin is not much better than it looks.} Every seen pull of bin $J$ has mean $\le m_J$, so
\[
\UCB_J\ \le\ m_J+2w(n_{i,J}).
\]

\emph{Combine.} The pulled arm has mean $\ge m_J-\omega$, and $\UCB_J\ge\UCB_{j^*}$. Therefore
\[
r_{i,t}\ \le\ m_{j^*}-(m_J-\omega)\ \le\ (\UCB_{j^*}+\omega)-(\UCB_J-2w(n_{i,J}))+\omega\ \le\ 2w(n_{i,J})+2\omega.\qedhere
\]
\end{proof}

\subsection{Stale counts}

This is the only place where the graph enters the regret bound.

\begin{lemma}[Unseen pulls]\label{lem:unseen}
Let $t$ be a Phase-2 round, $J$ a bin and $l\neq i$ an agent. The number of pulls by $l$ into $J$ in rounds
$e(t)<t'\le t-1$ is at most $\gamma\max(G_J(e(t)),1)+2L-1$.
\end{lemma}
\begin{proof}
Let $t$ be in slot $m+1$, so $\beta_m<t\le\beta_{m+1}$ and $e(t)=\max\{e\in\cE:e\le\beta_{m-1}\}$. Split the rounds
$(e(t),t-1]$ into $(e(t),\beta_{m-1}]$ and $(\beta_{m-1},t-1]$.

\emph{Second part.} It has $t-1-\beta_{m-1}\le\beta_{m+1}-1-\beta_{m-1}=2L-1$ rounds, and $l$ pulls once per round.

\emph{First part.} If $e(t)=\beta_{m-1}$, it is empty. Otherwise $\beta_{m-1}\notin\cE$, so no agent raised an alarm at
$\beta_{m-1}$. In particular, \eqref{eq:trigger} fails for $l$ and $J$ at $\beta_{m-1}$:
\[
O_{l,J}(\beta_{m-1})-O_{l,J}(\epsilon_{m-1})\ \le\ \gamma\max\big(G_J(\epsilon'_{m-1}),1\big).
\]
Here $\epsilon_{m-1}=\max\{e\in\cE:e\le\beta_{m-2}\}$. Because $\beta_{m-1}\notin\cE$, this equals
$\max\{e\in\cE:e\le\beta_{m-1}\}=e(t)$. Also $\epsilon'_{m-1}\le e(t)$, and $G_J$ is non-decreasing, so
$G_J(\epsilon'_{m-1})\le G_J(e(t))$. Hence $l$ made at most $\gamma\max(G_J(e(t)),1)$ pulls into $J$ in
$(e(t),\beta_{m-1}]$.
\end{proof}

\begin{lemma}[Stale counts]\label{lem:stale}
If the pull of agent $i$ at round $t$ is the $k$-th Phase-2 pull into bin $J$ (in lexicographic order), then
\[
n_{i,J}(t)\ \ge\ \frac{k-c_{N,\Lambda}}{1+N\gamma},\qquad c_{N,\Lambda}:=1+N(\gamma+2L)=1+N(\gamma+4\Lambda).
\]
\end{lemma}
\begin{proof}
Write $n=n_{i,J}(t)$ and $G=G_J(e(t))\le n$. The $k-1$ earlier pulls into $J$ are of three kinds:
\begin{itemize}[nosep]
\item seen by agent $i$: exactly $n$ of them (Lemma~\ref{lem:view});
\item by another agent $l\neq i$ in rounds $(e(t),t-1]$: at most $(N-1)\big(\gamma\max(G,1)+2L-1\big)$ by Lemma~\ref{lem:unseen};
\item by another agent in round $t$ itself: at most $N-1$.
\end{itemize}
So $k-1\le n+(N-1)\gamma\max(G,1)+(N-1)\,2L$.

\emph{Case $G\ge1$.} Then $\max(G,1)=G\le n$, so $k-1\le n(1+N\gamma)+2NL$, and
$n\ge(k-1-2NL)/(1+N\gamma)\ge(k-c_{N,\Lambda})/(1+N\gamma)$.

\emph{Case $G=0$.} Then $k-1\le n+N(\gamma+2L)$, so $n\ge k-c_{N,\Lambda}$. If $k\le c_{N,\Lambda}$ the claim is trivial.
Otherwise $n\ge k-c_{N,\Lambda}\ge(k-c_{N,\Lambda})/(1+N\gamma)$.
\end{proof}

\subsection{The regret theorem}

\begin{theorem}[Regret of Dec-ZoomSIB-Relay]\label{thm:regret}
Let Assumptions~\ref{ass} and~\ref{ass:delivery} hold, and let $\delta=1/(NT)$. Suppose $NT_0$ satisfies the
sample-size condition of Lemma~\ref{lem:stein}. Then for any $\Delta>0$ and $\gamma>0$,
\begin{align*}
\E R_T(N)\ \le\ {}&\underbrace{2L_fN(T_0+\Lambda)}_{\text{Phase 1}}
+\underbrace{4\sqrt2\,c_\delta\sqrt{(1+N\gamma)\Nb NT}}_{\text{UCB}}
+\underbrace{2L_{f'}\Delta NT}_{\text{discretisation}}\\
&+\underbrace{4L_{f'}\Xmax\,\eps_0(NT_0)\,NT}_{\text{projection}}
+\underbrace{2L_f\Nb\,(c_{N,\Lambda}+1)}_{\text{delay and staleness}}
+6L_f .
\end{align*}
\end{theorem}
\begin{proof}
\emph{Bad event.} $\Prob(E^c)\le3\delta=3/(NT)$, and the regret is always at most $2L_fNT$. This contributes at most $6L_f$.

\emph{Phase 1.} There are $N T_1=N(T_0+\Lambda)$ uniform pulls, each costing at most $2L_f$.

\emph{Phase 2 on $E$.} On $E_\theta$, $\eps_{\rm p}\le\Xmax\eps_0(NT_0)$. Fix a bin $j$ and let $n_j$ be its final number of
Phase-2 pulls. For the $k$-th of them, Lemmas~\ref{lem:perpull} and~\ref{lem:stale} give
\[
r^{(k)}\ \le\ \min\Big\{2L_f,\ 2c_\delta\sqrt{\tfrac{2(1+N\gamma)}{k-c_{N,\Lambda}}}+2\omega\Big\},
\]
where the second term is used only when $k>c_{N,\Lambda}$. Use $2L_f$ for $k\le c_{N,\Lambda}+1$; at most
$c_{N,\Lambda}+1$ such terms. For $k>c_{N,\Lambda}+1$, the numbers $k-c_{N,\Lambda}$ are at least $1$ and spaced by $1$,
so $\sum_k(k-c_{N,\Lambda})^{-1/2}\le\sum_{u=1}^{n_j}u^{-1/2}\le2\sqrt{n_j}$. Hence
\[
\sum_{k=1}^{n_j}r^{(k)}\ \le\ 2L_f(c_{N,\Lambda}+1)+4\sqrt2\,c_\delta\sqrt{(1+N\gamma)\,n_j}+2\omega\,n_j .
\]
Sum over the $\Nb$ bins. Cauchy--Schwarz gives $\sum_j\sqrt{n_j}\le\sqrt{\Nb\sum_jn_j}\le\sqrt{\Nb NT}$. Finally,
$2\omega NT=2L_{f'}\Delta NT+4L_{f'}\eps_{\rm p}NT$.
\end{proof}

\begin{corollary}[Rate]\label{cor:rate}
Take $\Delta=(NT)^{-1/3}$, $\gamma=\gamma_0/N$ for a constant $\gamma_0$, and
\[
NT_0=\Big\lceil\Big(L_{f'}\Xmax\tfrac{4C_\theta d}{\mu^*}\sqrt{\log(2dNT)}\,\tfrac{NT}{L_f}\Big)^{2/3}\Big\rceil
\]
(the choice of the server theorem), at least as large as Lemma~\ref{lem:stein} requires. Then
\[
\E R_T(N)=\tilde O\Big((d\,NT)^{2/3}+N\Lambda\,(NT)^{1/3}\Big).
\]
If $N\Lambda\le(NT)^{1/3}$, this matches Theorem~\ref{thm:lower} in $N$ and $T$ up to logarithms.
\end{corollary}
\begin{proof}
$\Nb\le2\Xmax(NT)^{1/3}+1=\tilde O((NT)^{1/3})$, and $c_\delta=\tilde O(1)$, so with $N\gamma=\gamma_0$ the UCB term
is $\tilde O((NT)^{2/3})$. The discretisation term is $2L_{f'}(NT)^{2/3}$. The choice of $T_0$ equalises
$2L_fNT_0$ and the projection term $\propto d\,NT/\sqrt{NT_0}$, and makes both $\tilde O((dNT)^{2/3})$. The Phase-1
delay $2L_fN\Lambda$ is dominated by the last term. The last term is
$2L_f\Nb(2+\gamma_0+4N\Lambda)=\tilde O(N\Lambda(NT)^{1/3})$.
\end{proof}

\begin{remark}[What changed from the sketch]
The earlier sketch (\texttt{server.tex}, Theorem~16) bounded the confidence width of an agent's view, which is a
union of $N$ per-agent prefixes, and lost a factor $\sqrt N$. Here the view is a common prefix plus one own block
(Lemma~\ref{lem:view}), so Lemma~\ref{lem:conf} is exactly the server lemma. The price is that agents act on
slightly older data: the delay constant is $4\Lambda$ instead of $\Lambda$. This is the ``delayed but consistent
statistics'' idea used by Dubey \& Pentland (2020) and by DLUCB.
\end{remark}

\begin{remark}[The trigger threshold]
The factor $\sqrt{1+N\gamma}$ appears for the same reason as with a server: between two epochs, each of the $N$
agents may add $\gamma\,G_J$ pulls to bin $J$. E30 measured this cost for the server protocol only. Whether
$\gamma=0.5$ is also nearly free under flooding has not been measured yet.
\end{remark}

%======================================================================
\section{Communication}\label{sec:comm}
%======================================================================

We count \emph{scalars}: integers or reals sent over one edge in one direction. This is the measure used by
Wang et al.\ (2020).

\begin{lemma}[Number of epochs]\label{lem:epochs}
For each bin $j$, let $A_j$ be the set of epochs at which some agent's alarm condition \eqref{eq:trigger} holds for bin
$j$. Then
\[
|A_j|\ \le\ 2\big(2+\log_{1+\gamma}(NT)\big),\qquad
|\cE|\ \le\ 1+2\Nb\big(2+\log_{1+\gamma}(NT)\big).
\]
\end{lemma}
\begin{proof}
List $A_j=\{a_1<a_2<\dots\}$ and let $g_p=G_j(a_p)$. Each $a_p$ is a slot boundary, so $a_{p+2}\ge a_p+2L$.

\emph{Growth over two steps.} Let $a_{p+2}=\beta_m$. The alarm at $\beta_m$ is by some agent $l$, with
$\delta_{l,j}(\beta_m)>\gamma\max(G_j(\epsilon'_m),1)$. The pulls counted in $\delta_{l,j}(\beta_m)$ happen in rounds
$(\epsilon_m,\beta_m]$, after $\epsilon'_m$. So
\[
g_{p+2}=G_j(\beta_m)\ \ge\ G_j(\epsilon'_m)+\delta_{l,j}(\beta_m)\ >\ G_j(\epsilon'_m)+\gamma\max\big(G_j(\epsilon'_m),1\big).
\]
Also $a_p\le\beta_m-2L=\beta_{m-2}$ and $a_p\in\cE$, so $a_p\le\epsilon'_m$ and $G_j(\epsilon'_m)\ge g_p$. Two cases follow:
\begin{itemize}[nosep]
\item if $g_p\ge1$, then $g_{p+2}>(1+\gamma)\,G_j(\epsilon'_m)\ge(1+\gamma)\,g_p$;
\item if $g_p=0$, then $g_{p+2}\ge\delta_{l,j}\ge1$.
\end{itemize}

\emph{Counting.} Consider the subsequence $g_1,g_3,g_5,\dots$. After at most one step it is $\ge1$. After that it grows by
a factor $>1+\gamma$ per step, and it never exceeds $NT$. So it has at most $2+\log_{1+\gamma}(NT)$ elements. The same
holds for $g_2,g_4,\dots$.

Finally, every epoch other than $\beta_0$ belongs to some $A_j$, and only bins that are actually pulled can trigger.
\end{proof}

\begin{theorem}[Communication]\label{thm:comm}
The total number of scalars sent by Dec-ZoomSIB-Relay is at most
\[
\underbrace{2|E|\,N\,(d+1)}_{\text{Phase 1}}\ +\ \underbrace{2|E|\,|\cE|\,\big(1+N(3\Nb+1)\big)}_{\text{Phase 2}}
\ =\ O\big(|E|\,N\,\Nb^2\log_{1+\gamma}(NT)+|E|\,N\,d\big).
\]
In particular it grows only logarithmically in $T$ for fixed $\Nb$. By contrast, any protocol in which every agent
sends its $\Nb$-bin table to its neighbours every round sends $\Theta(|E|\,\Nb\,T)$ scalars.
\end{theorem}
\begin{proof}
Under flooding, each message is forwarded by each agent at most once, to each neighbour. So one message crosses each
of the $2|E|$ directed arcs at most once.

\emph{Phase 1.} There are $N$ records, each with $d+1$ scalars.

\emph{Phase 2.} For each epoch, every agent sends at most one alarm over each arc: either its own alarm or the first
one it receives, never both. This is one scalar per arc. It also generates $N$ snapshots. Each snapshot holds at most $\Nb$ changed bins, at 3
scalars each (index, count increment, sum increment), plus one label. Multiply by $2|E|$ and by the number of epochs
from Lemma~\ref{lem:epochs}.
\end{proof}

\begin{remark}
Snapshots can instead be summed up a BFS spanning tree and broadcast back down. Each tree edge then carries at most
$3\Nb+1$ scalars per direction per epoch, which removes the factor $N|E|/(N-1)$:
$O\big((N\Nb+|E|)\,\Nb\log_{1+\gamma}(NT)\big)$. We keep flooding in the main statement because it is what the
robustness results of Section~\ref{sec:robust} need. The bound counts every pulled bin as changed in every snapshot, so
it is loose. In practice a snapshot contains only the bins the agent pulled since the last epoch.
\end{remark}

%======================================================================
\section{Robustness}\label{sec:robust}
%======================================================================

The proofs of Sections~\ref{sec:regret} and~\ref{sec:comm} use the graph only through Assumption~\ref{ass:delivery}.
So any communication model with a known delivery bound $\Lambda$ gets Theorems~\ref{thm:regret} and~\ref{thm:comm}
with that $\Lambda$. It suffices to compute $\Lambda$ in each case.

\begin{lemma}[Delivery time under weaker connectivity]\label{lem:delivery}
Assumption~\ref{ass:delivery} holds with the following $\Lambda$:
\begin{enumerate}[nosep,label=(\roman*)]
\item \emph{Directed graphs.} If $\cG$ is a static, strongly connected directed graph with directed diameter $D$,
then $\Lambda=D$.
\item \emph{Time-varying graphs.} Let the edge set $E_s$ change with the step $s$. Suppose that for every window of
$B$ consecutive steps, the union graph $\bigcup_{s\in\text{window}}E_s$ is connected with diameter at most $D_B$.
Use \emph{persistent flooding}: every agent re-sends each message it holds to its current neighbours in every step,
until $\Lambda$ steps after the message was created. Then $\Lambda=B(D_B+1)$.
\item \emph{Link delays.} Let $\cG$ have diameter $D$, and suppose every message sent over a link is delivered within
$\tau$ steps. Then $\Lambda=\tau D$.
\end{enumerate}
\end{lemma}
\begin{proof}
(i) This is the proof of Lemma~\ref{lem:flood}, with ``neighbour'' replaced by ``in-neighbour'' and distance by
directed distance.

(ii) Partial windows waste at most $B$ steps, so consider full windows after the one in which the message was sent.
Let $U_r$ be the set of agents that hold the message after $r$ full windows. Take an agent $u\in U_r$ and an edge
$\{u,v\}$ that appears at some step of window $r+1$. Under persistent flooding, $u$ sends the message across that edge
at that step, so $v\in U_{r+1}$. Hence $U_{r+1}$ contains $U_r$ and its neighbourhood in the union graph of window
$r+1$. That union graph has diameter at most $D_B$, so $U_{D_B}=V$.

(iii) The induction of Lemma~\ref{lem:flood}, with each hop taking at most $\tau$ steps.
\end{proof}

\begin{corollary}[Robustness]\label{cor:robust}
Under any model of Lemma~\ref{lem:delivery}, with slot length $L=2\Lambda$ for the corresponding $\Lambda$,
Theorem~\ref{thm:regret} and Corollary~\ref{cor:rate} hold verbatim. In particular, a $B$-connected time-varying graph
costs an additive $\tilde O\big(N B D_B(NT)^{1/3}\big)$ in regret. Records arrive late but are never wrong, so the
leading $(NT)^{2/3}$ term does not change. Theorem~\ref{thm:comm} holds verbatim in cases (i) and (iii). In case (ii)
it holds with an extra factor $\Lambda$, and with $|E|$ replaced by $\max_s|E_s|$, because each message is re-sent for
$\Lambda$ steps.
\end{corollary}
\begin{proof}
Lemma~\ref{lem:welldef} is the only place where delivery is used, and it uses only Assumption~\ref{ass:delivery}.
Everything after it depends only on the epoch structure. For communication, persistent flooding sends each message
over each present arc at most once per step, for $\Lambda$ steps.
\end{proof}

\begin{remark}
Random link failures, such as each link failing independently with probability $q$ per step (one of the settings
of E18), give a random delivery time. With persistent flooding, the corollary then applies on the event that every
message used is delivered within a chosen $\Lambda$. A union bound over messages and over the hops of a shortest path
shows that this event has probability at least $1-1/(NT)$ when $\Lambda\gtrsim(D+\log(NT))/(1-q)$. The complement adds
a term of order $L_f$ to the bound.
\end{remark}

%======================================================================
\section{Agreement}\label{sec:agree}
%======================================================================

So far every agent froze the same direction $\hat\theta_0$. Approximate Phase-1 methods such as consensus or
push-sum leave each agent $k$ with its own direction $\hat\theta_k$. Bins are then computed with different
projections, but all tables are still added up bin by bin.

\begin{theorem}[Cost of disagreement]\label{thm:agree}
Run Phase 2 of Dec-ZoomSIB-Relay with agent $k$ binning by its own $\hat\theta_k$, $\|\hat\theta_k\|_1=1$, on the same grid.
Let $\eta=\max_k\|\hat\theta_k-\theta^*\|_1$. On the event $\{\eta\le\bar\eta\}$, Theorem~\ref{thm:regret} holds with
the projection term replaced by
\[
4L_{f'}\Xmax\,\bar\eta\,NT .
\]
In particular, if $\hat\theta$ is the exact pooled direction and $\max_k\|\hat\theta_k-\hat\theta\|_1\le\eta_{\rm agree}$,
then the extra regret relative to exact pooling is at most $4L_{f'}\Xmax\,\eta_{\rm agree}\,NT$.
It is zero when Phase~1 pools exactly.
\end{theorem}
\begin{proof}
Lemmas~\ref{lem:geom}, \ref{lem:conf}, \ref{lem:perpull} and~\ref{lem:stale} were stated for arbitrary directions
$\hat\theta_k$. The directions enter only through $\eps_{\rm p}=\Xmax\max_k\|\hat\theta_k-\theta^*\|_1$:
\begin{itemize}[nosep]
\item Lemma~\ref{lem:geom}(ii) puts the true index of a bin-$j$ pull by \emph{any} agent in the same widened interval
$\tilde I_j$. So pooling the bin-$j$ statistics of different agents is still pooling pulls whose means lie in
$[m_j-\omega,m_j]$.
\item Lemma~\ref{lem:conf} concerns only noise, which does not depend on the directions.
\item Lemma~\ref{lem:stale} concerns only counts.
\end{itemize}
The proof of Theorem~\ref{thm:regret} then goes through with $\eps_{\rm p}\le\Xmax\bar\eta$. For the second statement,
use the triangle inequality $\|\hat\theta_k-\theta^*\|_1\le\|\hat\theta-\theta^*\|_1+\|\hat\theta_k-\hat\theta\|_1$ and
compare with the exact-pooling bound.
\end{proof}

\begin{corollary}[An agreement pass is enough]\label{cor:agreepass}
Suppose Phase~1 uses any method that leaves agent $k$ with a Stein sum $b_k'$ (an estimate of the pooled sum), followed
by one flooding pass in which every agent floods its own \emph{local} sum $b_k$. Then all agents end with the exact
pooled direction after $\Lambda$ more rounds, and Theorem~\ref{thm:regret} holds with $T_0+\Lambda$ replaced by
$T_0+2\Lambda$ in the Phase-1 term.
\end{corollary}
\begin{proof}
After the flooding pass every agent holds all $N$ local sums (Assumption~\ref{ass:delivery}), so all agents compute
$\sum_kb_k/\|\sum_kb_k\|_1$. Lemma~\ref{lem:exact} applies, and so does Theorem~\ref{thm:agree} with $\eta_{\rm agree}=0$.
The extra $\Lambda$ rounds of uniform play cost at most $2L_fN\Lambda$.
\end{proof}

\begin{remark}[Link to E28]
Theorem~\ref{thm:agree} is linear in the disagreement and has no threshold. E28 (fig28, panel D) shows regret rising
steeply once disagreement exceeds about $10^{-2}$. That is consistent, since the bound only becomes comparable to the
other terms once $\Xmax\eta_{\rm agree}NT$ is of order $(NT)^{2/3}$. E28b confirms Corollary~\ref{cor:agreepass}: one flooding pass
removes the whole penalty of one-step consensus (from $\times1.25$ to $\times1.08$ on the ring, and from $\times1.34$ to
$\times1.09$ on the path).
\end{remark}

%======================================================================
\section{Summary}
%======================================================================

\begin{center}
\begin{tabular}{@{}lll@{}}
\toprule
Result & Statement & Evidence \\ \midrule
Thm~\ref{thm:lower} & $\Omega((NT)^{2/3})$ for any graph and protocol & E22 \\
Thm~\ref{thm:regret}, Cor~\ref{cor:rate} & $\tilde O((dNT)^{2/3}+N\Lambda(NT)^{1/3})$ & E19, E23 (code version)\\
Thm~\ref{thm:comm} & $O(|E|N\Nb^2\log_{1+\gamma}NT)$ scalars & E15, E20, E23 \\
Cor~\ref{cor:robust} & time-varying, directed, lossy graphs: replace $D$ by $\Lambda$ & E18 \\
Thm~\ref{thm:agree} & disagreement $\eta$ costs $\le4L_{f'}\Xmax\eta NT$ & E28 \\
\bottomrule
\end{tabular}
\end{center}

\paragraph{Open.} (a) The additive $N\Lambda\Nb$ term: E23 measures a cost 60--150$\times$ smaller that grows like
$D^{0.65}$. (b) A lower bound that depends on the graph. (c) The trigger cost under flooding (E31). (d) An
experiment with the epoch protocol of Section~\ref{sec:alg}, to confirm it behaves like the code version.

\end{document}
